% Generated draft: requires mathematical and source review.
\documentclass[11pt]{article}
\usepackage[margin=25mm]{geometry}
\usepackage{fontspec}
\setmainfont{Latin Modern Roman}
\usepackage{amsmath,amssymb,mathtools,unicode-math}
\setmathfont{Latin Modern Math}
\usepackage{enumitem,needspace,xurl,hyperref}
\setlength{\parindent}{0pt}
\setlength{\parskip}{5pt}
\setlength{\emergencystretch}{3em}
\newcommand{\N}{\mathbb{N}}
\newcommand{\Z}{\mathbb{Z}}
\newcommand{\Q}{\mathbb{Q}}
\newcommand{\R}{\mathbb{R}}
\newcommand{\C}{\mathbb{C}}
\newcommand{\F}{\mathbb{F}}
\newcommand{\A}{\mathbb{A}}
\newcommand{\PP}{\mathbb{P}}
\newcommand{\E}{\mathbb{E}}
\newcommand{\eps}{\varepsilon}
\newcommand{\dd}{\,\mathrm{d}}
\begin{document}
\section*{UnsolvedMath ID: 30006312}
\textbf{Source problem code:} catalogue entry
\subsection{Definitions and mathematical statement}
A global field \(F\) is a finite extension of \(\Q\) or a finite extension of \(\F_q(t)\), where \(q\) is a prime power. Let \(S\) be a finite set of places of \(F\); if \(F\) is a number field, require \(S\) to contain every archimedean place. For each nonarchimedean place \(v\), write \(v:F^\times\to\Z\) for its normalized additive valuation and set \(v(0)=+\infty\). The ring of \(S\)-integers is
\[
\mathcal O_{F,S}=\{x\in F: v(x)\geq 0\text{ for every nonarchimedean place }v\notin S\}.
\]
The supplied record imposes no nonemptiness condition on \(S\) in the function-field case.

\textbf{Source question.} For a global field \(F\) and a finite set \(S\) satisfying these conditions, how does the intended arithmetic--topology correspondence change when its base ring is replaced by \(\mathcal O_{F,S}\)?

\textbf{Unresolved scope.} This is an extension question, not a specified theorem or equivalence. The supplied record does not identify the original base ring, define the arithmetic--topology correspondence, or specify the objects, constructions, or properties to be compared after replacement. Consequently, it does not determine a unique precise conjecture beyond the question above. The original excerpt uses \(F\) for the field but \(K\) in the number-field qualification; the supplied corrected statement interprets both as \(F\), and this draft follows that interpretation.
\subsection{Short English statement}
For a global field \(F\), the ring \(\mathcal O_{F,S}\) permits denominators at the nonarchimedean places in a finite set \(S\), which must also include all archimedean places in the number-field case. The question asks how the intended arithmetic--topology correspondence extends to this ring. The supplied excerpt leaves the correspondence and the desired conclusions unspecified.
\subsection{Sources}
\begin{itemize}
\item Local UnsolvedMath registry, permanent ID \(30006312\), title: \emph{Arithmetic Topology over Rings of S-Integers}; source problem code \texttt{catalogue entry}, in the supplied \emph{Oberwolfach Reports Open Problems} collection.
\item \emph{Arbeitsgemeinschaft: Relative Langlands Duality} (2025), Oberwolfach Reports, pp.~811--882, DOI \texttt{10.4171/owr/2025/18}: \url{https://doi.org/10.4171/owr/2025/18}. The supplied original excerpt includes a report-title marker followed by page number 829. No individual proposer is named; the registry supplies proposed year 2025.
\item Supplied literature source, labeled ``Oberwolfach report'': \url{https://doi.org/10.4171/owr/2025/18}. The literature assessment dated \texttt{2026-08-22} labels the question open and reports that no general theorem answering the extension to rings of \(S\)-integers was located. This is supplied dated curation, not an independently verified literature assessment; the background expressly requests independent verification of time-sensitive status and cited claims.
\item The registry assigns statement status \texttt{corrected\_verified} and reports that the official report was checked, the question isolated from surrounding exposition, and its notation repaired for self-contained publication. These are supplied provenance claims, not verification performed for this draft.
\end{itemize}
\end{document}
